\begin{abstract}

When an LLM writes code to satisfy a multi-step specification, does it know \emph{what it intends to do next} before it starts writing, or is it improvising token by token? We show that code LLMs form an explicit representation of their next instrumental sub-goal \emph{before} they emit a single token of its implementation, and that this representation is \textbf{decodable}, and \textbf{portable}. First, we localize the intent of the next instrumental goal to a compact set of attention heads, which we call the \goalheads. Using what the \goalheads write to the residual stream we can get a faithful readout of the upcoming sub-task well before the LM realizes this in tokens. We also find that this encoding is \emph{abstract}: transplanting it steers the model to pursue the intended sub-goal in a fresh context, in a different programming language. Furthermore, find that the same \goalheads give faithful readouts of the LM's next-step in writing the code where the specifications are not concretely defined in the prompt. The same attention heads encodings give faithful readout and remain causal in non-code settings; such as few-shot relation mappings settings and plannings tasks. Our results give a concrete and testable evidence that LLM's in their latents encode abstract representations of intermediate goals, while pursuing long-horizon tasks.

\end{abstract}
